% =====================================================================
%  MELBA submission — Journal of Machine Learning for Biomedical Imaging
%  Ported from the TUBITAK (elektr.cls) version in ../v2-improved/.
%
%  MELBA is SINGLE-blind: author names and affiliations belong ON this PDF.
%  Do NOT modify melba.cls — "using any other template will result in a
%  desk reject" (melba-journal.org, verified 2026-07-25).
%
%  Compile:  pdflatex manuscript && bibtex manuscript && pdflatex x2
% =====================================================================
\documentclass{melba}

% melba.cls already provides: amsmath, natbib, graphicx, booktabs, multirow,
% xcolor, hyperref, lineno, microtype. Only additions go here.
\usepackage{amsfonts}

% --- Float placement -------------------------------------------------------
% In two-column mode LaTeX may only set a starred (full-width) float at the top
% of a page, and with 11 tables and 5 figures competing it deferred every figure
% past the bibliography: Figures 2-5 were being typeset on pages 22-25, after
% the references. dblfloatfix repairs double-column float ordering and permits
% bottom placement; the loosened fractions stop TeX from refusing a page that
% is mostly float. Without this the figures do not appear near the text that
% discusses them.
\usepackage{dblfloatfix}
\renewcommand{\topfraction}{0.90}
\renewcommand{\bottomfraction}{0.75}
\renewcommand{\textfraction}{0.08}
\renewcommand{\floatpagefraction}{0.70}
\renewcommand{\dblfloatpagefraction}{0.70}
\renewcommand{\dbltopfraction}{0.90}
\setcounter{topnumber}{3}
\setcounter{bottomnumber}{2}
\setcounter{totalnumber}{5}
\setcounter{dbltopnumber}{3}

% Allow long URLs in the reference list to break across lines.
\def\UrlBreaks{\do\/\do\-\do\.\do\_\do\?\do\=\do\&\do\#}
\Urlmuskip=0mu plus 1mu\relax

% Two-column measure is narrow and this text is full of long hyphenated
% compounds ("source-grouped", "single-dataset", "burn-severity"). Allow TeX
% a last-resort stretch so no line is set overfull, and teach it the few
% compounds it cannot break on its own.
\emergencystretch=3em
\hyphenation{Al-moa-ma-ri da-ta-set seg-men-ta-tion clas-si-fi-er
             per-cep-tual re-pro-duc-i-bil-i-ty gen-er-al-i-sa-tion}

% --- Header/footer fields. Filled by the publishing editor on acceptance;
%     placeholders are correct for an initial submission. ---
\melbaid{YYYY:NNN}
\doi{}
\melbaauthors{Altayeb, Alraddadi and Alrehaili}
\email{altayeb.ray@gmail.com}
\volume{}
\firstpageno{1}
\melbayear{2026}
\datesubmitted{}
\datepublished{}

\ShortHeadings{Six evaluation choices that changed our conclusions}{Altayeb, Alraddadi and Alrehaili}

% Title history, kept because the reasons are substantive.
%  1. "A comparative study of deep learning approaches for automated burn injury
%     segmentation and severity classification" -- enrolled the paper in the exact
%     genre its own Section 2 criticises, and signalled none of its contribution.
%  2. "A leak-free, externally validated benchmark ..." -- dropped because
%     Section 4.11 reports that our own pipeline leaked, so a title claiming
%     "leak-free" contradicted the paper's most important self-correction.
%  3. "What a second look changed: five protocol audits ..." -- dropped as too
%     vague to survive a contents page: it names neither the domain nor the
%     finding, and "protocol audit" is not standard terminology.
% The present title puts the domain first, as a medical-imaging readership
% expects, and states the contribution plainly in the second clause. It claims
% only what the paper demonstrates.
\title{Deep learning for burn severity assessment: six evaluation choices
that changed our conclusions}

% Affiliation convention: you list the institution where the work was
% PERFORMED, not where you happen to be today. This study originated as the
% authors' undergraduate graduation project at the Islamic University of
% Madinah, so that is the correct affiliation and removing it would make the
% paper less accurate, not more. The authors have since graduated; that fact
% is stated plainly in the Acknowledgements instead of being hidden.
\author{
	\name Ryan Altayeb\aff{1}\orcid{0009-0006-0031-9531},
	\name Abdulrahman Alraddadi\aff{1}\orcid{0009-0000-2325-5063},
	\name Mohannad Alrehaili\aff{1}\orcid{0009-0008-2491-185X}
}
\affiliations{%
	\num 1 \addr Faculty of Computer and Information Systems, Islamic University of Madinah, Madinah, Saudi Arabia
}

\abstract{%
Burn-depth assessment guides clinical management but is subjective, and deep learning is an attractive aid. The burn literature, however, reports mostly single-dataset accuracy, with little external validation or scrutiny of data hygiene. We study a deployed two-stage system --- a YOLOv8x-seg localiser feeding a masked Swin classifier, served through a mobile application and framed as a research artifact rather than a validated device --- and build a benchmark around it on a source-grouped split of 1,370 photographs, with 11 architectures, two external test sets and a ten-seed replication. We then re-examine six routine evaluation choices under an independent second analysis.
%
Splitting the augmented files at random left 80.5 percent of the masked test set sharing a source photograph with training while the unmasked split leaked none, manufacturing a masking benefit of 2.07 percentage points on the eight architectures common to both analyses; grouped by source photograph, and using ground-truth masks so the figure bounds what any real segmenter could deliver, it is 0.55 points (95 percent confidence interval $-0.37$ to $+1.47$). Grouping by source identifier, the standard correction, still left 2 of 205 test images perceptually identical to a training image filed under another identifier. Three training seeds put the internal head-to-head margin at $+3.74$ points with an interval about three times narrower than ten seeds support, and detected no external margin at all; ten seeds give $+2.63$ ($+0.29$ to $+4.98$) internally and $+2.79$ ($+1.09$ to $+4.49$) externally, and of all 120 three-seed subsets of those same ten runs only 15 detect the latter. A single-dataset subgroup test produced a large skin-tone association ($p=0.0025$, surviving Bonferroni correction) that reversed and vanished on a better-powered external cell. Our own corrected split reached only the classifiers: 179 of 205 internal test images (87.3 percent) lay in the segmentation models' training or validation folds, and retraining them cost the standalone model 12.7 points internally, a figure confounded with the 21 percent smaller training set the correction implies. Finally, the internal margin holds on plain accuracy but not on balanced accuracy ($+1.45$, $p=0.17$); only the external margin survives both.
%
A plain classifier was the most accurate single system throughout, so segmentation-first is justified by localisation rather than accuracy; on an independent clinical database the model did not transfer at all, and on a perceptual-hash-cleaned web-sourced set it under-graded 28.3 percent of eligible images against 19.3 internally. Two of the six answers vanished, one reversed, one was overstated while its companion was concealed, one exposed a leak we had introduced ourselves, and one held only under the metric we happened to have chosen.
%
Code, weights, per-image predictions and a script recomputing every published number: \url{https://github.com/Just-Ryan/burn-severity-benchmark}.}

\keywords{Burn severity assessment, deep learning, data leakage, external validation, image segmentation, reproducibility}

\begin{document}

\twocolumn[\maketitle]

\section{Introduction}
\label{Int}
Burns are common and serious injuries. About 11 million people need medical attention for burns each year, and the World Health Organization estimates that burns cause roughly 180{,}000 deaths annually, most of them in low- and middle-income countries where specialist care is scarce \citep{who,peck}. Survivors often face long hospital stays, scarring, and disability \citep{brusselaers}. Deciding how to treat a burn depends first on assessing its depth, which separates injuries that heal with conservative care from those that need surgery \citep{jackson}. Burns are graded as first degree (epidermis only), second degree (partial thickness, into the dermis), or third degree (full thickness) \citep{heimbach}.

Depth assessment is difficult. Even experienced burn surgeons assess burn depth accurately only 64 to 76 percent of the time \citep{jaskille}, and it is affected by lighting, wound evolution, blistering, and skin pigmentation \citep{monstrey}. Instrumented alternatives such as infrared thermography reach high accuracy but are expensive and rarely available \citep{dang}. An automated, reproducible assessment aid could therefore be useful, particularly where specialists are absent.

Deep learning has reached expert-level performance on several medical imaging tasks \citep{esteva,litjens}, and it has been applied to burns for two coupled subproblems, segmenting the burn region and classifying its severity \citep{boissin}. Progress is real, but the burn literature has three recurring gaps. First, most studies train and report a single architecture, which makes it hard to know whether a reported number reflects the method or the particular network. Second, few studies report confidence intervals, significance tests, or evaluation on data from a source other than the training set. Third, and rarely examined, is data hygiene: whether the split procedure itself inflates the reported numbers. These gaps matter because high in-distribution accuracy has repeatedly failed to predict external performance in medical imaging, and because the evaluation choices that produce such numbers are themselves under active scrutiny: leakage has recently been shown to inflate prediction performance across a range of mechanisms in neuroimaging \citep{rosenblatt2024}, and community guidance on metric selection for biomedical image analysis was substantially revised in 2024 \citep{metricsreloaded}. A recent five-year systematic review of the burn literature finds accuracies of 80 to 98 percent reported on single datasets, judges six of fourteen studies to carry a high risk of bias, faults the reliance on public datasets that do not represent the target population, and reports dataset details as frequently unstated \citep{fu2026review}; the evidence that region-of-interest masking helps at all is likewise mixed \citep{sourget}. Our central thesis, threaded throughout this paper, is that in-distribution accuracy is not evidence of clinical usefulness, whereas leak-free external evaluation is.

We make five contributions. \textbf{First}, and the reason the rest exists, we measure how far six routine evaluation choices move a published answer, each against a corrected baseline (Table~\ref{tab:audit}). Splitting an augmented dataset at the file level manufactured a masking benefit of 2.07 percentage points on the eight architectures common to both analyses, which falls to 0.55 (95 percent confidence interval $-0.37$ to $+1.47$) once the split is grouped by source photograph. Grouping by source identifier, the standard correction for exactly that failure, still left 2 of 205 internal test images (1.0 percent) perceptually identical to a training image filed under a different identifier, one pair carrying conflicting labels. A three-seed protocol overstated our internal head-to-head margin, narrowed its interval about threefold, and concealed our external margin entirely; across all 120 three-seed subsets of a ten-seed replication only 15 detect the external effect that ten seeds resolve at $p=0.005$, and the estimated standard deviation ranges from 0.00 to 3.25 percentage points. A single-dataset subgroup test produced a large, Bonferroni-surviving skin-tone association that reversed direction and vanished on a better-powered external cell. And having built that source-grouped split for the classifiers, we left the segmentation models on the collection's original split, so 179 of the 205 internal test images (87.3 percent) sat in their training or validation folds; every internal number passing through a YOLO model is consequently optimistic, in a direction that makes our own headline margin a lower bound rather than an overstatement.

\textbf{Second}, we build the leak-free benchmark on which those measurements were made: a source-grouped split of 1,370 photographs, 11 ImageNet-pretrained architectures under four input conditions, one internal and two external test sets, Wilson intervals and paired significance tests kept separate by estimand, and a ten-seed replication of the head-to-head. \textbf{Third}, we report two distinct results on whether segmentation-first helps. Masking the image to the burn confers no benefit resolvable at this sample size, $+0.55$ percentage points with a 95 percent interval of $-0.37$ to $+1.47$; because that ablation used ground-truth masks it is an upper bound, so a better segmenter could not rescue it. Across training runs the plain classifier does not merely tie the pipeline but beats it, by 2.63 percentage points internally (95 percent interval $+0.29$ to $+4.98$) and 2.79 externally ($+1.09$ to $+4.49$), ahead in nine of ten external runs. Two qualifications travel with that: on any single set of images the difference remains unresolved, and on balanced rather than plain accuracy only the external margin survives. We keep the estimands and the summaries apart throughout. \textbf{Fourth}, we report that the system does not transfer to independent clinical images and fails in the dangerous direction, under-grading 28.3 percent of the images that could be under-graded on a perceptual-hash-cleaned external set against 19.3 percent internally. \textbf{Fifth}, we describe the deployed system that motivated all of this, a YOLOv8x-seg localiser feeding a masked Swin classifier served through a Flutter and Flask mobile application, and we state plainly what it is not. It is not a state-of-the-art architecture, and it is not a clinically validated tool. It is a research and demonstration artifact and, as our external results show, it is not ready for any clinical role.

We did not anticipate the first of those contributions when we began. Six times in this study a conventional analysis and a second, independent look gave different answers, and in every case the conventional protocol was the one we had run first. The masking benefit vanished once the split was grouped by source photograph, and the standard identifier-based grouping that achieved this still left duplicates behind. The skin-tone association reversed. The internal head-to-head margin shrank and its interval widened about threefold, while its external counterpart, the most secure result in the paper, was invisible until we stopped using three seeds. And our own leak-free split covered only the stage of the pipeline we had thought to rebuild. We therefore treat the paper's methodological findings as its primary contribution and the burn-severity results as the setting in which they were obtained.

The remainder of the paper is organised as follows. Section~\ref{Rel} reviews related work. Section~\ref{Mat} describes the dataset, the leak-free split and its correction, the deployed system, the benchmark design, the training protocol, and the evaluation and statistical methods. Section~\ref{Res} reports the results in the order segmentation quality; the leakage artifact and the masking null; the head-to-head comparison as originally run; what a ten-seed replication changed; oracle against predicted masks; external failure; what any three of those ten seeds would have reported; the skin-tone probe; the segmentation-split leak we found in our own pipeline; and efficiency. Section~\ref{Dis} interprets these results and states the limitations, and Section~\ref{Con} concludes.

\begin{table*}[t]
\centering
\caption{The six routine evaluation choices audited in this paper, what each appeared to show, and what the corrected analysis shows. No new data were collected between the two columns. Rows 1, 3, 5 and 6 recompute on data already in hand; row 4 corrects by comparison against an independent set, which is itself the audited failure there; row 2 is a content-based audit of the same split. Here and in the tables that follow, pp abbreviates percentage points.}
\small
\begin{tabular}{|p{0.16\textwidth}|p{0.28\textwidth}|p{0.38\textwidth}|c|}
\hline
\textbf{Routine choice} & \textbf{Apparent result} & \textbf{Result after correction} & \textbf{Section} \\
\hline
Split the augmented files at random
& Masking helps: $+2.07$ pp on the 8 architectures common to both analyses ($+3.06$ across all 21 originally trained, 20 of 21 positive)
& $+0.55$ pp, 95\% CI $[-0.37,+1.47]$, Wilcoxon $p=0.32$, 13 of 24 architecture-seed pairs positive; ground-truth masks, so an upper bound
& \ref{Res:leak}, \ref{Res:mask} \\
\hline
Group by source identifier and assume the folds are independent
& The folds are independent
& 2 of 205 internal test images (1.0\%) are perceptually identical to a training image filed under a different identifier, one pair with conflicting labels; on a second public collection, 118 images and 56 of 199 sources (28\%) were duplicates of our training data
& \ref{Mat:leak}, \ref{Mat:eval} \\
\hline
Use three training seeds
& Internal $+3.74$ $[+3.04,+4.44]$, $p=0.002$; external $+3.24$ $[-3.83,+10.31]$, $p=0.19$
& Ten seeds: internal $+2.63$ $[+0.29,+4.98]$, $p=0.032$; external $+2.79$ $[+1.09,+4.49]$, $p=0.005$, ahead in 9 of 10. Across all 120 three-seed subsets of those ten runs: SD 0.00 to 3.25 pp, external $p$ from 0.0005 to 0.86, only 15 of 120 reach $p<0.05$
& \ref{Res:stat}, \ref{Res:seeds} \\
\hline
Test a subgroup on one dataset
& Errors fall on darker-appearing skin: median ITA 19.2 against 38.4, $p=0.0025$, surviving Bonferroni correction
& External replication gives $p=0.35$ with the effect reversed (rank-biserial $+0.10$) on a larger cell, 139 images against 81
& \ref{Res:ita} \\
\hline
Summarise accuracy the usual way
& The classifier beats the pipeline internally: $+2.63$ pp $[+0.29,+4.98]$, $p=0.032$, over ten seeds
& On \emph{balanced} accuracy over the same ten runs the internal margin is $+1.45$ pp $[-0.74,+3.65]$, $p=0.17$, ahead in 6 of 10, and is not established. The external margin survives both summaries ($+1.96$ $[+0.26,+3.65]$, $p=0.028$)
& \ref{Res:stat} \\
\hline
Correct the split for one stage and assume the pipeline is clean
& The head-to-head is leak-free, because the classifier split is grouped by source photograph
& The segmentation models were trained earlier on the collection's original split: 179 of 205 internal test images (87.3\%) lie in their training or validation folds. Retraining both models on the source-grouped split drops the standalone model from 90.7\% to \textbf{78.0\%} on the same images (exact McNemar $p=4\times10^{-5}$). That 12.7 pp is confounded with a 21\% smaller training set; a 26-image control neither model saw puts 7.7 pp of it outside leakage, with a wide interval. All external numbers, and every result computed from ground-truth masks, are unaffected (verified, zero overlap)
& \ref{Res:segleak} \\
\hline
\end{tabular}
\label{tab:audit}
\end{table*}

\section{Related work}
\label{Rel}
\textbf{Automated burn assessment and the high single-dataset pattern.} Computer-aided burn assessment predates deep learning: early systems on the Seville BIP\_US database combined hand-crafted colour and texture features with classical classifiers to separate burn depths \citep{acha}. Deep learning has since become dominant and is typically applied to two coupled subproblems, delineating the burn region and grading its severity \citep{boissin}. Recent work spans convolutional classifiers \citep{suha}, attention-augmented residual networks \citep{yadav2023}, YOLO-family detectors adapted to burn severity \citep{ferdinand,yildiz}, and transformer-based frameworks that perform segmentation and severity assessment jointly \citep{nadeem}. Several report high single-dataset accuracy, for example 95.6 percent with a transfer-learning classifier \citep{suha}, and a recent five-year systematic review reports single-dataset accuracies of 80 to 98 percent while judging six of fourteen studies to be at high risk of bias, largely because the public datasets used do not represent the target population \citep{fu2026review}. Which severity class is hardest is not settled in this literature: \citet{nadeem} report second degree as among their best-performing classes and attribute their weakest classes to under-representation, whereas second degree is consistently the hardest class in our own results.

\textbf{Comparative studies on a single split.} Comparative burn studies exist \citep{rahman,chang2023}, but they evaluate several backbones on a single split and rarely report confidence intervals, significance tests, or external evaluation. Without these, apparent improvements may not be reliable and out-of-distribution behaviour is unknown. Our work targets exactly these gaps: we compare approaches under one fixed protocol, report Wilson confidence intervals, apply paired significance tests, and add explicit external tests.

\textbf{Segment-then-classify as an established design, and whether masking helps.} Cascading a segmentation stage into a downstream classifier is well established in medical imaging: dual-stage frameworks localise a region of interest and then classify it in breast ultrasound \citep{bruno} and in dermoscopy \citep{almasni2020}, and the same cascade appears in burns for wound measurement and surface-area estimation \citep{chang2021}. This grounds our two-stage design as a standard template, not a novelty. The usual motivation is that masking to the lesion suppresses background shortcuts and focuses the classifier on relevant tissue. The evidence that this helps is mixed: models trained on full images can match or exceed models restricted to the region of interest because networks exploit background and spurious cues, so masking is not a universally beneficial preprocessing step \citep{sourget}. This is consistent with our central negative result.

\textbf{Mobile and point-of-care deployment and its validation gap.} Because our contribution is a deployed system, the most relevant comparators are fielded tools. The most clinically validated is the DeepView multispectral device, whose artificial-intelligence ensemble predicted burn healing potential with 100 percent specificity in a feasibility study \citep{thatcher2023}, albeit with dedicated multispectral hardware rather than a consumer camera. Smartphone-only wound tools are more common but less validated: a deep-learning pressure-injury staging app reported 63.2 percent accuracy on its 144-image validation set \citep{lau2022}, and a review of mobile ulcer-segmentation apps concluded that most disclose neither their training data nor adequate clinical validation \citep{griffa2024}. Against this backdrop, a consumer-phone burn system such as ours, like the mobile burn app of \cite{yildiz}, is valuable precisely when it is reported honestly, including where it underperforms.

\textbf{Data leakage, external validation, and the central theme.} Leakage, information about the test set contaminating training, is a long-recognised cause of inflated, irreproducible results \citep{kaufman}, and it has been identified as a systemic driver of the reproducibility crisis across machine-learning-based science \citep{kapoor}. In medical imaging, a systematic review that screened 2{,}212 studies and adjudicated 62 found none of them of potential clinical use, with leakage, biased curation, and inadequate reporting among the leading causes \citep{roberts}, and community responses include consensus reporting standards such as the REFORMS checklist \citep{reforms}. A large body of evidence shows that in-distribution accuracy is a poor predictor of external accuracy: a pneumonia network that scored an internal area under the curve of 0.931 fell to 0.815 at an external hospital and was found to exploit site-specific artifacts \citep{zech2018}; a systematic review of externally validated radiology models found external performance routinely lower and frequently unreported \citep{yu2022}; and reliable in-distribution magnetic-resonance models degraded sharply on out-of-distribution cohorts \citep{martensson2020}. A concrete and easily reproduced form of leakage in image work is splitting after augmentation, so that augmented copies of one source image land in both folds. Quantified case studies of this kind exist, and they report inflations larger than ours. Splitting brain MRI at the slice rather than the subject level inflated accuracy by between 29 and 55 percent across four cohorts \citep{yagis2021}, and an equivalent failure in retinal and breast OCT inflated accuracy by 5 to 30 percent \citep{tampu2022}. The picture is not uniform: in connectome-based prediction, some forms of leakage inflate performance substantially while others have minor or even negative effects \citep{rosenblatt2024}. What these cases share is that leakage moves the absolute score of a single model. Closer to our case, confound removal, applied as a corrective, has itself been shown to introduce leakage into one arm of a comparison and so to manufacture an apparent benefit for that arm \citep{hamdan2023}. Our study places burn classification within this pattern and closes on the theme: in-distribution accuracy is not evidence of clinical usefulness; leak-free external evaluation is.

\textbf{How many training runs a conclusion needs.} That a benchmark result depends on more than the reported point estimate is established. Comparing two methods from single runs produces high rates of both false positives and false negatives, and the variance contributed by perturbing the data split exceeds that contributed by weight initialisation alone \citep{bouthillier2021} --- an ordering our own results reproduce, since our split procedure moved the masking answer further than our seed did. Even where the overall spread across seeds is modest, individual outlier seeds that perform markedly above or below the average are easy to find \citep{picard2021}. In medical imaging specifically, evaluation noise has been shown to exceed the gap between the winning model and the best 10 percent of entries in five of eight challenges examined \citep{varoquaux2022}, and biomedical image analysis rankings are unstable to methodological choices such as the metric variant or the removal of a single test case \citep{maierhein2018}. What we add is not the observation that seeds matter, which belongs to this literature, but a measured instance in which the conventional three-run protocol changed which of two conclusions the evidence supported, in both directions at once (Section~\ref{Res:stat}), together with the full distribution of answers the same data would have yielded (Section~\ref{Res:seeds}).

\section{Materials and methods}
\label{Mat}
\subsection{Dataset}
\label{Mat:data}
The primary dataset was the publicly available ``skin burn wound classification'' dataset on Roboflow Universe (version 31) \citep{roboflow}, released under a Creative Commons Attribution 4.0 licence and exported in COCO segmentation format. It comprises \textbf{1{,}370 unique burn photographs} in three severity classes; the export augments these into about 3{,}424 annotated files with distinct file names (about 2.5 copies per source). Table~\ref{tab:data} summarises the dataset, the leak-free split, and the two external sets. We map the numeric source labels to first, second, and third degree. Class shares are approximately 42 percent first degree, 31 percent second, and 27 percent third. The labels are community-provided and were not independently verified by clinicians.

\begin{table*}[t]
\centering
\caption{Dataset, leak-free split, and external sets. The naive file-level split (documented only to explain the leakage artifact) shared a source photograph between test and training for 80.5 percent of masked test images and none of the unmasked test images.}
\small
\begin{tabular}{|l|c|}
\hline
Item & Value \\
\hline
\multicolumn{2}{|l|}{\textbf{Primary set (Roboflow ``skin burn wound classification'', version 31, CC BY 4.0)}} \\
\hline
Unique source photographs & 1{,}370 \\
Augmented files & $\approx$ 3{,}424 \\
Split (by source, stratified by class, seed 42) & 70 / 15 / 15 \\
Validation / test (deduplicated to one image per source) & 206 / 205 \\
Internal test distribution (first / second / third) & 86 / 63 / 56 \\
Naive-split source overlap into training (masked / unmasked) & 80.5\% / 0\% \\
\hline
\multicolumn{2}{|l|}{\textbf{External sets}} \\
\hline
BIP\_US clinical database (Seville; second-vs-third probe) & 94 \\
Clean external set (second Roboflow collection, pHash-verified) & 319 \\
\quad distribution (first / second / third) & 139 / 157 / 23 \\
\quad images removed as perceptual-hash identical to training & 118 \\
\quad sources contributing a removed image & 56 / 199 (28\%) \\
\quad sources eliminated entirely & 24 / 199 \\
\hline
\end{tabular}
\label{tab:data}
\end{table*}

From the split we derive a segmentation dataset (polygon annotations, either a single ``burn'' class or three severity classes) and a classification dataset in four input forms: unmasked full images; masked images, in which the ground-truth polygon annotation is rasterised to a binary mask and multiplied into the image so that non-burn pixels are zero; images cropped to the annotation bounding box; and masked and cropped. We emphasise that the masked form uses ground-truth annotations rather than predicted masks, so any benchmark result on it is an oracle upper bound, a distinction we return to when interpreting the pipeline.

\subsection{Leak-free split and the leakage correction}
\label{Mat:leak}
The naive split follows common practice and partitions the augmented image files at random. Because the source dataset augments each photograph into several copies with distinct file names, a file-level split scatters copies of one photograph across training and test. In the masked classification set built this way, re-split 70/15/15 after flattening, \textbf{417 of its 518 test images (80.5 percent) shared a source photograph with training}. The unmasked set was never re-split: it retained the collection's own 90/6/4 partition, whose 143-image test fold leaked none. The two arms of that original comparison therefore differed in split ratio and in test-set size as well as in leak rate, which we state because it makes the confound larger than a single-split description would suggest, not smaller. This asymmetry, together with memorisation of near-duplicate training images, manufactured an apparent masking benefit of 3.06 percentage points across 20 of 21 architectures. That figure and the leak-free figure below come from different architecture pools, so we also report the matched contrast: restricted to the eight architectures common to both analyses, the naive split gives $+2.07$ percentage points (7 of 8 positive) against $+0.55$ on the source-grouped split. The matched contrast, $2.07 \rightarrow 0.55$, is the smaller of the two and is the one we rely on; we give the 21-architecture figure because it is what the original experiment produced. We report the leaky result only to document the artifact, not as performance, and we never use it as a headline result.

The correction is a source-grouped split. We partition the 1{,}370 source photographs 70/15/15, stratified by class and seeded at 42, keep all augmented copies only in the training fold, and deduplicate validation and test to one image per source (206 and 205 images). All downstream leak-free numbers use this split. The internal test set thus contains \textbf{205} images, distributed 86 first degree, 63 second degree, and 56 third degree.

Source grouping removes leakage between augmented copies of the same file, but it cannot remove what the file names do not encode. Auditing the corrected split with the same perceptual-hash procedure we later apply to the external set (Section~\ref{Mat:eval}), we find \textbf{2 of the 205 test images (1.0 percent) are perceptually identical to a training image carrying a different source identifier}, that is, the same photograph was contributed to the collection twice under different names. One of the two pairs is additionally annotated with conflicting labels in the two folds. We report this rather than suppress it, for three reasons. First, the effect is bounded, though not as harmlessly as we first wrote. Every reported accuracy moves by at most 1.0 percentage point, which is within the cross-seed standard deviation. But the paired ten-seed internal difference is itself only $+2.63$ points, and if both images were counted as errors for the classifier and correct for the pipeline on every seed --- the worst case, and an implausible one, but a $4/205$ swing rather than a $2/205$ one --- that difference falls to $+0.68$ with a 95 percent interval of $-1.66$ to $+3.03$ and $p=0.53$. Even the milder scenario in which only the classifier loses the two images leaves $+1.66$, interval $-0.69$ to $+4.00$. Neither excludes zero. The external comparison, which carries the paper's argument, is unaffected because these two images are internal. We record this because an earlier draft of this paragraph claimed the duplicates ``reverse no comparison in this paper'', and for the internal paired comparison under its worst case that was not true. Second, it is a concrete instance of the paper's own argument, namely that leakage is easy to introduce and hard to see, and that a single grouping heuristic is not a proof of independence. Third, it yields a practical recommendation we now follow ourselves: \emph{grouping by identifier should be verified with a content-based check}, because identifiers describe provenance only as well as the upstream collection was curated.

\subsection{Deployed two-stage system}
\label{Mat:system}
The deployed system processes an image in two stages (Figure~\ref{fig:system}). A single-class YOLOv8x-seg localiser predicts a burn mask; the mask is binarised and multiplied into the image so that the background is black; the masked image is passed to a Swin transformer classifier (Swin-Tiny), which outputs the severity grade. The system is served through a Flask backend and a Flutter mobile front-end. The mobile demonstration currently ships a Swin-Small classifier, whereas all pipeline numbers reported in this paper use the benchmarked Swin-Tiny weights. We stress that the application is a research and demonstration artifact, not a clinically validated device, and Section~\ref{Res:ext} shows that it does not transfer to independent clinical images.

\begin{figure*}[t]
\centering
\includegraphics[width=0.86\textwidth]{fig_system.pdf}
\caption{The deployed two-stage system. A mobile application sends an image to a Flask backend, where a YOLOv8x-seg localiser predicts a burn mask, the mask is applied so that the background is black, and a Swin classifier grades the masked region as first, second, or third degree. The system is a research and demonstration artifact, not a clinically validated device.}
\label{fig:system}
\end{figure*}

\subsection{Benchmark design}
\label{Mat:bench}
We run two benchmarks that answer two distinct questions.

Benchmark 1 asks whether focusing the classifier on the burn helps. We train 11 ImageNet-pretrained architectures, a representative subset of the 21 used in the naive-split experiment chosen to span convolutional and transformer families within the compute budget, on the leak-free split: ConvNeXt-Tiny, ConvNeXt-Large, DenseNet201, EfficientNet-B0, MobileNetV3-Large, ResNet50, Swin-Tiny, Swin-Small, Swin-Base, Swin-Large, and ViT-Base \citep{convnext,densenet,efficientnet,mobilenet,resnet,swin,vit}. Coverage of the four input conditions of Section~\ref{Mat:data} is not uniform, and we state it explicitly: all 11 architectures were trained under the unmasked and masked conditions, whereas the two cropping conditions were run on 8 of the 11, omitting ConvNeXt-Tiny, Swin-Base, and Swin-Large for compute reasons. Because the masked condition uses ground-truth masks, its result is an oracle upper bound on any masking benefit that a real segmenter could achieve.

Benchmark 2 asks whether segmentation-first helps end to end, through a head-to-head on the shared 205-image internal test set and on the external sets. The two arms are not symmetric, and we state the asymmetry rather than claim fairness. The standalone classifier is evaluated in its best Benchmark 1 condition, ConvNeXt-Large on unmasked images, whereas the pipeline arm is fixed to the deployed configuration, Swin-Tiny on masked images, because the deployed system is what this paper is about. This favours the classifier arm, and the comparison should be read as the best available plain classifier against the deployed pipeline configuration, not as the best of each family. The conclusion does not rest on that choice alone: Benchmark 1 independently bounds the \emph{average} masking benefit at about $+1.5$ percentage points even with ground-truth masks. That is a bound on the pooled mean across architectures, not on any single one, and Section~\ref{Res:mask} shows the per-architecture effect is not stable; it therefore constrains, without excluding, the possibility that some untested masked configuration does better. The robust pipeline runs the localiser at a low confidence threshold of 0.05, masks the image, and classifies it with Swin-Tiny, falling back to the full image when the localiser returns no region, which matches the deployed application. The all-in-one YOLOv8-seg model is evaluated at the same low confidence threshold of 0.05 for fairness. We report accuracy and, on the imbalanced external sets, balanced accuracy.

For robustness, the Benchmark 2 classifiers were trained and evaluated over seeds 0, 1, and 2 and are reported as mean and sample standard deviation (denominator $n-1$), while the Benchmark 1 masking ablation was confirmed over seeds 0, 1, and 42 for the masked-versus-unmasked comparison on 8 of the architectures; the cropping conditions were evaluated at seed 42 only. These are separate experiments, run at different times, and we report each seed set as it stands rather than merging them.

\subsection{Training protocol}
\label{Mat:train}
Classifiers used ImageNet-pretrained backbones \citep{imagenet} from the timm library \citep{timm}, each with the classification head resized to three outputs and trained under one fixed recipe: the AdamW optimiser, class-weighted cross-entropy using inverse class frequencies, a cosine schedule with early stopping, mixed precision, and input resized to $224\times224$. Training-time augmentation (horizontal flip, colour jitter, small rotation) was enabled and evaluation-time augmentation disabled, with an identical transform for every architecture so that differences reflect the network rather than the setup.

Both YOLOv8x-seg models (the single-class localiser and the three-class standalone model) used the same configuration: 100 epochs at $640\times640$ input, batch size 8, the AdamW optimiser at an initial learning rate of $10^{-3}$, seed 42, and early stopping with patience 20, initialised from the public \texttt{yolov8x-seg.pt} weights \citep{ultralytics}.

\textbf{An important asymmetry between the two stages, which we discovered late and report in full.} The source-grouped split of Section~\ref{Mat:leak} was constructed for the \emph{classification} datasets. The segmentation models were trained earlier, on the collection's original distribution-supplied split, and we did not rebuild them when we rebuilt the classifier split. The two splits do not agree. Auditing them against each other, \textbf{179 of the 205 internal classification test images (87.3 percent) fall in the segmentation models' training or validation folds}: 145 (70.7 percent) in training and 34 in validation, leaving only 26 images that those models had never seen. Section~\ref{Res:segleak} reports what this does and does not affect. In short, every internal number that passes through a YOLO model is optimistic; the segmentation quality figures, the entire Benchmark 1 masking ablation, the oracle-mask results and every external number are unaffected, and we verified rather than assumed each of those.

The Benchmark 1 masking ablation was run on a Kaggle NVIDIA P100 (PyTorch 2.5.1, CUDA 12.1) \citep{pytorch}. The Benchmark 2 segmentation training, the head-to-head classifier retraining, the external evaluation, and the efficiency measurements were run on a Google Colab NVIDIA A100 (PyTorch 2.11, CUDA 12.8). The retrained classifier accuracies are about two percentage points lower than the earlier seed-42 Kaggle run. Hardware, library version, recipe and seed set all differ between the two, so we cannot attribute that gap to any one of them and do not try; we report it so that the two sets of numbers are not mistaken for replicates of each other.

\subsection{Evaluation, robust pipeline, and statistics}
\label{Mat:eval}
For a like-for-like comparison, the three-class standalone YOLOv8-seg detections are reduced to a per-image majority-degree label, directly comparable to the classifier and the pipeline on the same 205 images. Images for which the localiser returns no region are never dropped from the denominator. Under the strict pipeline they are counted as errors; under the robust pipeline they fall back to the full frame and are scored on that, which is what the deployed application does. To make the effect of this choice transparent, we report both a strict pipeline (default confidence 0.25, no-detection counted as an error) and the robust pipeline (confidence 0.05 with full-image fallback) as an ablation.

We report accuracy with Wilson 95 percent confidence intervals \citep{wilson} and compare systems on the same images with the paired exact McNemar test \citep{mcnemar}. Two distinct questions require two distinct tests, and we keep them separate throughout. Whether one system beats another \emph{on a particular test set} is an image-level question, answered by McNemar on paired predictions. Whether one system beats another \emph{on average across training runs} is a run-level question, answered by a paired test over per-seed accuracies. Overlapping error bars are not used as evidence of absence, since the comparisons are paired.

The masking effect was tested with the Wilcoxon signed-rank test over architecture-seed pairs, and we report it as an effect size with a confidence interval rather than against a single detectability threshold. The difference between transformer and convolutional architectures was tested with a one-sided Mann-Whitney test (the directional hypothesis being that transformers benefit more); the two-sided value is also reported in Section~\ref{Res:mask} for transparency. Following the honesty rule of this study, we prefer interval estimates to accept/reject statements wherever the data allow.

Two external protocols probe generalisation. The first uses the BIP\_US clinical database of the Biomedical Image Processing Group at the University of Seville (94 photographs) \citep{acha,rostami}, which predates and is independent of the web-sourced training data. BIP\_US labels depth as superficial dermal, deep dermal, or full thickness; because the mapping onto our three classes is not exact, we report both plausible mappings of deep dermal (to second degree and to third degree), giving a second-versus-third-degree probe with no first-degree cases. The second external set is a clean set of 319 images that we built from a second public Roboflow burn collection by removing every image that overlapped our training data. We screened for overlap with a perceptual-hash method: pHash and dHash computed with the \texttt{imagehash} library \citep{zauner}, compared by Hamming distance against all primary training images and made horizontal-flip-aware. An image was treated as a duplicate when either hash matched a training image exactly, at Hamming distance 0, in any of the flip orientations tested; after removal we verified that no retained image lay within a Hamming distance of 10 of any training hash. The collection contained 1{,}371 images from 199 source photographs. The check removed 118 images that were perceptual-hash identical to training; those images came from 56 of the 199 sources (28 percent). Removal was at the image level, not the source level, so a source that contributed one contaminated image could retain its other copies, and it is the verified Hamming-distance-10 margin above, not source-level exclusion, that rules out residual contamination among the retained images. After also dropping the fourth-degree class and remapping the remainder to first, second, and third degree, 319 images remain, distributed 139, 157, and 23, drawn from \textbf{175 distinct source photographs} (mean 1.8 images per source, maximum 4); 24 of the original 199 sources are gone entirely. We state this arithmetic explicitly because 199 minus 56 is not 175, and a reader is right to check. they are external to the training data, but they are not 319 statistically independent observations, so an interval computed on $N=319$ is too narrow. Rather than only note this, we quantify it: alongside every external Wilson interval we report a cluster bootstrap that resamples the 175 \emph{source photographs} with replacement (20{,}000 resamples, percentile method), which is the interval we rely on. Clustering widens the external accuracy intervals by roughly a fifth, from 9.6 to 11.7 percentage points for the classifier and from 10.0 to 11.9 for the pipeline. We retain the image-level point estimates for comparability with the internal test set. One consequence is worth stating because it is a defence rather than a concession: the ten-seed paired differences of Section~\ref{Res:stat} are computed over \emph{training runs}, not over images, so source clustering does not widen those intervals. That defence is about the interval, not about generalisation: any claim beyond these particular 319 images still inherits the test set's limited source diversity. The perceptual-hash check is reported as a method, not an afterthought, because it reinforces the leakage and reproducibility theme.

\subsection{Skin-tone probe}
\label{Mat:ita}
Because any imaging aid intended for clinical use should be checked for systematic dependence on skin tone, we added an exploratory probe. No ground-truth phototype labels exist for either dataset, so we estimated \emph{apparent} skin tone from the images with the Individual Typology Angle, $\mathrm{ITA}=\arctan((L^{*}-50)/b^{*})\cdot 180/\pi$, computed in CIE $L^{*}a^{*}b^{*}$ under a D65 white point and binned into the six conventional bands of \citet{chardon1991} (very light, light, intermediate, tan, brown, dark).

ITA must be measured on unburned skin, since a burn changes colour drastically. We therefore excluded the burn region from every image before sampling: on the internal set from the ground-truth mask, recovered from the masked and unmasked copies of the same photograph, and on the external set by rasterising the polygon annotations. Both exclusions were dilated by 6 pixels so that the burn margin could not contaminate the sample. Within the remaining region we selected skin pixels by the conjunction of the standard RGB \citep{kovac2003} and YCbCr skin rules and took the median ITA over those pixels. Images yielding fewer than 500 skin pixels were recorded as indeterminate and excluded rather than guessed at; this removed 11 of the 205 internal images and none of the external images, leaving 194 and 319.

We report accuracy per band with Wilson intervals, but we test the association on the \emph{continuous} ITA with a Mann-Whitney test comparing the ITA of correctly and incorrectly classified images, so that no conclusion depends on where the band edges happen to fall.

The limits of this measurement are severe and we state them before reporting any result. ITA computed from uncontrolled photographs of unknown provenance is a proxy for apparent skin tone \emph{in the image}, not a measurement of a patient's constitutive phototype: white balance, illumination, camera processing and perilesional erythema all shift it. The probe is exploratory and hypothesis-generating. It is not a fairness audit, and we do not present it as one.

\section{Results}
\label{Res}
\subsection{Segmentation quality}
\label{Res:seg}
On the held-out test set the single-class localiser, which is the pipeline front-end, reached a mask mean average precision at intersection-over-union 0.50 (mask mAP50) of \textbf{0.726} and a mask mAP50-95 of 0.417, while the three-class standalone model reached a mask mAP50 of \textbf{0.603} and a mask mAP50-95 of 0.349 (Table~\ref{tab:seg}). The single-class model scores higher because it need not separate severity classes. Second degree is the hardest class to segment, consistent with the classifiers' weakest per-class performance on partial-thickness burns. Segmentation is therefore only moderate, a point that matters for the pipeline below.

\begin{table*}[t]
\centering
\caption{Segmentation quality on the held-out test set (Benchmark 2, YOLOv8x-seg). Both models used the identical configuration (100 epochs, $640\times640$, batch 8, AdamW at $10^{-3}$, seed 42, patience 20). Second degree is the hardest class.}
\small
\begin{tabular}{|l|c|c|c|}
\hline
Model & Classes & Test mask mAP50 & Test mask mAP50-95 \\
\hline
Single-class localiser (pipeline front-end) & 1 & 0.726 & 0.417 \\
Standalone YOLOv8-seg & 3 & 0.603 & 0.349 \\
\hline
\end{tabular}
\label{tab:seg}
\end{table*}

\subsection{The apparent masking benefit was leakage}
\label{Res:leak}
Under the naive file-level split, masking appeared to raise accuracy for 20 of 21 architectures by a mean of 3.06 percentage points, and by 2.07 points across the eight architectures that also enter the pooled leak-free estimate (7 of 8 positive), which is the matched figure the leak-free result should be compared against (Figure~\ref{fig:leak}). This benefit does not survive scrutiny: it is a data-leakage artifact, driven by the 80.5 percent versus 0 percent source overlap between the masked and unmasked test sets (Table~\ref{tab:data}), where the masked accuracies were inflated by memorisation of near-duplicate training images.

What distinguishes this from the quantified leakage case studies of Section~\ref{Rel} is not its size, which is smaller than theirs, but its structure. In those studies leakage inflates the absolute score of a model. Here the leak rate was \emph{confounded with the experimental condition under comparison}, 80.5 percent in the masked arm against none in the unmasked arm, so leakage did not inflate a number but manufactured a comparative conclusion and, from it, a design recommendation, doing so consistently across 20 of 21 architectures. The nearest precedent we have found is confound-leakage, where a corrective preprocessing step introduces leakage into one arm of a comparison and so produces an apparent benefit for that arm \citep{hamdan2023}. Our own mechanism is narrower than confound-leakage and we should not conflate the two: masking introduces no leakage of itself, and what differed between our arms was the split protocol, one arm re-partitioned at file level and the other left on the collection's own partition (Section~\ref{Mat:leak}). Split-after-augmentation leakage is itself documented \citep{yagis2021,tampu2022}. What we have not found previously documented is the combination, an asymmetric split protocol that leaks into only one arm of a controlled comparison, with the per-arm leak rate measured on both sides; that combination, rather than any single ingredient, is what turned a null into a design recommendation. The transferable lesson is narrow and practical: when a leak rate differs between the arms being compared, consistency across architectures is evidence \emph{for} the confound rather than against it, and is therefore not the reassurance it appears to be.

We document the leaky result only to explain the artifact. Every conclusion below rests on the source-grouped split, with the one exception recorded in Section~\ref{Res:segleak}, where the segmentation stage was found to have kept the original partition.

\begin{figure*}[t]
\centering
\includegraphics[width=0.92\textwidth]{fig_leakage.pdf}
\caption{The data-leakage artifact at the centre of this study. (a) The source collection augments each photograph into about 2.5 near-duplicate copies with distinct file names, so partitioning the augmented files at random places copies of one photograph in both the training and the test fold. (b) The consequence, averaged over architectures: under the naive file-level split the masked condition appears to beat the unmasked one by 3.06 percentage points in 20 of 21 architectures, an artifact of that shared-photograph path. Under the source-grouped split the same comparison gives 1.51 percentage points at seed 42 and only 0.55 percentage points pooled over seeds 0, 1 and 42, with a 95 percent confidence interval spanning zero. Panel (b) plots the 21-architecture figure because that is what the original experiment produced; restricted to the eight architectures common to both analyses the naive split gives 2.07 percentage points, so the matched before-and-after contrast is $2.07 \rightarrow 0.55$. Either way, the effect that motivated the original design does not survive a leak-free split.}
\label{fig:leak}
\end{figure*}

\subsection{On a leak-free split, masking gives no reliable benefit}
\label{Res:mask}
On the leak-free split the masking benefit disappears (Table~\ref{tab:mask}). For seed 42 alone, background masking raised accuracy by 1.5 percentage points overall, concentrated in the transformer architectures (3.1 percentage points, in all five of five) with roughly no change for the convolutional networks. That pattern did not replicate. Adding the fresh seeds 0 and 1 to seed 42, over the 24 architecture-seed pairs available in all three runs, collapses the effect to \textbf{$+0.55$ percentage points with a 95 percent confidence interval of $-0.37$ to $+1.47$} (Wilcoxon signed-rank $p=0.32$; 13 of 24 pairs improved). The interval is the informative quantity: it is consistent with no effect, and it rules out any benefit larger than about 1.5 percentage points. The apparent transformer advantage also failed to replicate (one-sided Mann-Whitney $p=0.33$; two-sided $p=0.66$), with transformers at $+0.77$ and convolutional networks at $+0.33$ percentage points pooled. Cropping to the burn produced essentially no change at seed 42, and masking with cropping likewise; the two cropping conditions were run on 8 of the 11 architectures and were not repeated across seeds, so we treat them as a weaker, single-seed check. We read this symmetrically: not that masking is exactly zero, but that any benefit it confers is smaller than the run-to-run variation of the training procedure itself, and that the seed-42 signal was noise. This matches the wider finding that region-of-interest masking is not reliably beneficial \citep{sourget}.

\begin{table*}[t]
\centering
\caption{Masking ablation (Benchmark 1) on the leak-free split. Each condition is compared to unmasked full images on the same architectures. The masked condition uses ground-truth polygon masks, so it is an oracle upper bound. Changes are means in percentage points; the architecture count for each condition is given explicitly.}
\small
\begin{tabular}{|l|c|c|c|c|}
\hline
Condition & Architectures & Seed 42 & Seeds 0, 1, 42 & Significance \\
\hline
Masked (background zeroed) & 11 & $+1.5$ & $+0.55$ & Wilcoxon $p=0.32$; CI $[-0.37,+1.47]$ \\
Cropped to bounding box & 8 & $\approx +0.0$ & not run & not significant \\
Masked and cropped & 8 & $\approx +0.0$ & not run & not significant \\
\hline
\end{tabular}

\vspace{2mm}
\parbox{0.95\textwidth}{\footnotesize The seed-42 sweep covered all 11 architectures under the masked and unmasked conditions, but only 8 of the 11 under the two cropping conditions (ConvNeXt-Tiny, Swin-Base, and Swin-Large were not run with cropping, for compute reasons). The multi-seed confirmation over seeds 0, 1, and 42 was run for the masked condition only, across 8 architectures and 24 architecture-seed pairs; the cropping conditions were evaluated at seed 42 only. The seed-42 masking effect was concentrated in transformers ($+3.1$, five of five) with roughly no change for convolutional networks ($+0.2$); the transformer-versus-convolutional difference was not significant (one-sided Mann-Whitney $p=0.33$, two-sided $p=0.66$). The pooled masking effect is $+0.55$ percentage points with a 95 percent confidence interval of $-0.37$ to $+1.47$; the interval, not a detectability threshold, is the quantity we rely on. For contrast, the naive file-level split had appeared to show $+3.06$ percentage points across 20 of 21 architectures, a leakage artifact. Restricted to the same eight architectures that enter the pooled leak-free estimate, the naive split gives $+2.07$ percentage points (7 of 8 positive), so the matched before-and-after contrast is $2.07 \rightarrow 0.55$ rather than $3.06 \rightarrow 0.55$.}
\label{tab:mask}
\end{table*}

\subsection{Head-to-head: a plain classifier is the most accurate single system}
\label{Res:h2h}
Table~\ref{tab:h2h} and Figure~\ref{fig:h2h} report the core comparison. Across seeds 0, 1, and 2, the standalone classifier (ConvNeXt-Large) reached \textbf{82.6 $\pm$ 0.7} percent internally and \textbf{76.6 $\pm$ 1.6} percent on the clean external set (balanced 80.1 $\pm$ 1.8), while the robust pipeline reached \textbf{78.9 $\pm$ 1.0} percent internally and \textbf{73.4 $\pm$ 2.5} percent externally (balanced 77.9 $\pm$ 3.4). The classifier won on both sets in all three seeds. The all-in-one YOLOv8-seg model, evaluated at the fair confidence threshold of 0.05 but trained with a single seed, reached 78.0 percent internally (balanced 77.7) and 66.8 percent externally (balanced 67.0). These are the figures for the model retrained on the source-grouped split (Section~\ref{Res:segleak}); we treat the point estimates cautiously because each is a single seed.

\begin{table*}[t]
\centering
\caption{Head-to-head accuracy (percent). Multi-seed rows are mean $\pm$ sample standard deviation (denominator $n-1$) over seeds 0, 1, and 2, with the Wilson 95 percent interval for the seed-42 run given beneath each. Balanced accuracy is given for the imbalanced external set. Reference rows below the rule isolate the oracle-mask upper bound and the threshold-sensitivity ablation and are not the deployed system.}
\small
\begin{tabular}{|l|c|c|}
\hline
System & Internal ($N=205$) & External ($N=319$) \\
\hline
Standalone classifier (ConvNeXt-Large) & $82.6 \pm 0.7$ & $76.6 \pm 1.6$ (bal.\ $80.1 \pm 1.8$) \\
\quad Wilson 95\% CI (seed 42) & [76.6, 87.0] & [69.2, 78.8] \\
\quad Source-clustered 95\% CI (seed 42) & n/a & \textbf{[68.3, 80.1]} \\
Robust pipeline (localiser to Swin-Tiny) & $78.9 \pm 1.0$ & $73.4 \pm 2.5$ (bal.\ $77.9 \pm 3.4$) \\
\quad Wilson 95\% CI (seed 42) & [75.6, 86.2] & [64.7, 74.7] \\
\quad Source-clustered 95\% CI (seed 42) & n/a & \textbf{[63.7, 75.7]} \\
Standalone YOLOv8-seg (single seed, conf.\ 0.05) & $\mathbf{78.0}$ (bal.\ $77.7$) & $\mathbf{66.8}$ (bal.\ $67.0$) \\
\quad same model before the split was corrected & \emph{90.7} (bal.\ \emph{90.6}) & \emph{48.9} (bal.\ \emph{50.3}) \\
\hline
Oracle Swin-Tiny (ground-truth masks) & $83.9$ & n/a \\
Strict pipeline (conf.\ 0.25, no-detection = error) & $79.0$ & $54.2$ \\
Robust pipeline (seed 42) & $81.5$ & $69.9$ \\
\hline
\end{tabular}

\vspace{2mm}
\parbox{0.95\textwidth}{\footnotesize The standalone YOLOv8-seg row is the model retrained on the source-grouped split (Section~\ref{Res:segleak}); the italicised row beneath it is the same architecture and recipe trained on the collection's original split, which had seen 87.3 percent of the internal test set. Both are given because the pair is informative: the contaminated model's internal-to-external drop is 41.8 points against the retrained model's 11.2, and that excess is the leakage signature. The internal gap between the two rows is not attributable to leakage alone, since the corrected partition also leaves 21 percent less training data (Section~\ref{Res:segleak}). The robust and strict pipeline rows still carry the \emph{original} localiser and their internal accuracies are therefore optimistic, which makes the classifier's internal margin over them a lower bound. The oracle-mask row and the classifier rows never involve a trained segmenter and are unaffected throughout.

\medskip
\textbf{Two questions, two tests.} \emph{Does the classifier beat the pipeline on these particular images?} Paired exact McNemar on the seed-42 models: internal 20 versus 18 discordant pairs, $p=0.87$ (a tie; seed 42 was the pipeline's most favourable run); external 54 versus 40, $p=0.18$. \emph{Does the classifier beat the pipeline on average across training runs?} The paired per-seed difference is $+3.74$ pp internally (95\% CI $+3.04$ to $+4.44$, paired $t$ $p=0.002$) and $+3.24$ pp externally (95\% CI $-3.83$ to $+10.31$, $p=0.19$), with the classifier ahead in three of three seeds on both sets. The internal advantage is therefore statistically supported at the level of training runs, while the single-run image-level comparison is not resolved at $N=205$. We note the paired tests use only three seeds and are correspondingly imprecise; Table~\ref{tab:paired} repeats them over ten.}
\label{tab:h2h}
\end{table*}

\begin{figure}[t]
\centering
\includegraphics[width=\linewidth]{head_to_head.pdf}
\caption{Head-to-head accuracy of the standalone classifier, the robust two-stage pipeline, and the all-in-one YOLOv8-seg model on the internal test set ($N=205$) and the clean external set ($N=319$). Error bars show the standard deviation across three seeds for the classifier and the pipeline; the standalone YOLOv8-seg model is single-seed and is shown without error bars. The classifier is highest on both sets and leads in every seed. Because the comparison is paired, the correct test is on the per-seed differences rather than on the overlap of these bars: that paired difference is $+3.74$ percentage points internally (95 percent CI $+3.04$ to $+4.44$) and $+3.24$ points externally (95 percent CI $-3.83$ to $+10.31$).}
\label{fig:h2h}
\end{figure}

\subsection{Three seeds overstated one margin and concealed the other}
\label{Res:stat}
The two questions of Section~\ref{Mat:eval} give different answers, and separating them matters. At the level of \emph{these particular images}, the advantage is not resolved: on the seed-42 models McNemar's exact test gave 20 versus 18 discordant pairs internally ($p=0.87$, effectively a tie, because seed 42 happened to be the pipeline's most favourable run) and 54 versus 40 externally ($p=0.18$).

At the level of \emph{training runs} the original three-seed protocol suggested a clear internal advantage and an unresolved external one. Neither conclusion survived more seeds in the form it was stated, and the way each failed is instructive enough to report in full. We repeated the entire head-to-head over \textbf{ten} seeds, training both arms under a fixed recipe and evaluating each on both test sets, with the pipeline arm using \emph{predicted} masks from the localiser rather than ground-truth masks, so that the arm being measured is the deployed system (Table~\ref{tab:paired}).

Over ten seeds the classifier leads the pipeline by \textbf{2.63} percentage points internally (95 percent CI $+0.29$ to $+4.98$; paired $t$ $p=0.032$, Wilcoxon $p=0.041$; ahead in 7 of 10 runs) and by \textbf{2.79} percentage points externally (95 percent CI $+1.09$ to $+4.49$; $p=0.005$, Wilcoxon $p=0.014$; ahead in \textbf{9 of 10} runs).

The choice of metric qualifies the internal result, and we report the qualification rather than the more convenient figure. Repeating both comparisons on \emph{balanced} accuracy, which is the appropriate summary when classes are unequal, the external margin survives at $+1.96$ points (95 percent CI $+0.26$ to $+3.65$, $p=0.028$, ahead in 8 of 10) but the internal margin does not: $+1.45$ points, 95 percent CI $-0.74$ to $+3.65$, $p=0.17$, ahead in only 6 of 10 runs. The internal advantage is therefore a property of plain accuracy on this class distribution and is not established once the classes are weighted equally, whereas the external advantage holds under both summaries. That a metric choice of this kind can decide a conclusion is the subject of a substantial recent literature on validation practice in biomedical imaging \citep{metricsreloaded}, and our case is a small instance of it. This is the sixth of the evaluation choices audited in this paper (Table~\ref{tab:audit}), and it points the same way as the other five: the external comparison is the one that carries the argument.

The contrast with three seeds is the point. Three seeds put the internal difference at $+3.74$ points with an interval of $+3.04$ to $+4.44$ and $p=0.002$: an interval about three times narrower than ten seeds support (1.40 against 4.69 percentage points wide), around an effect whose true magnitude is smaller. The same three seeds put the external difference at $+3.24$ points with an interval of $-3.83$ to $+10.31$ and $p=0.19$, wide enough to license no claim at all, when the effect is in fact the most reliable one in the study. A three-seed protocol thus overstated the weaker result and concealed the stronger one simultaneously.

We therefore state the head-to-head as follows. \emph{The standalone classifier outperforms the two-stage pipeline on both test sets on plain accuracy, by roughly two to three percentage points; on balanced accuracy only the external margin holds, and it is the external margin that is most reliable} where it holds in nine of ten training runs. This is a more modest internal claim and a considerably stronger external one than the original protocol implied, and the external claim carries the argument, since it concerns data the model was never fitted to.

This is the clearest evidence in the paper for its own thesis. A three-seed protocol, entirely ordinary in this literature, did not merely add noise: it produced a confidently narrow interval around an overstated effect while hiding a real one. Both errors survive peer review easily, because an interval computed from three runs is typographically indistinguishable from one computed from thirty.

\begin{table*}[t]
\centering
\caption{The head-to-head repeated over ten seeds, both arms trained under a fixed recipe and each evaluated on both test sets. The pipeline arm uses predicted masks from the localiser (confidence 0.05 with full-frame fallback), so it is the deployed system rather than the oracle-mask upper bound. Differences are paired by seed. The three-seed rows are the original protocol, shown for comparison.}
\small
\begin{tabular}{|l|c|c|c|c|}
\hline
 & Classifier & Pipeline & Paired difference (95\% CI) & $p$ \\
\hline
\multicolumn{5}{|l|}{\textbf{Internal test ($N=205$)}} \\
\hline
Three seeds (original) & $82.6 \pm 0.7$ & $78.9 \pm 1.0$ & $+3.74$ \ \ $[+3.04, +4.44]$ & 0.002 \\
\textbf{Ten seeds} & $83.6 \pm 1.7$ & $81.0 \pm 2.2$ & $\mathbf{+2.63}$ \ \ $\mathbf{[+0.29, +4.98]}$ & 0.032 \\
\hline
\multicolumn{5}{|l|}{\textbf{External test ($N=319$)}} \\
\hline
Three seeds (original) & $76.6 \pm 1.6$ & $73.4 \pm 2.5$ & $+3.24$ \ \ $[-3.83, +10.31]$ & 0.19 \\
\textbf{Ten seeds} & $77.8 \pm 2.3$ & $75.0 \pm 1.3$ & $\mathbf{+2.79}$ \ \ $\mathbf{[+1.09, +4.49]}$ & \textbf{0.005} \\
\hline
\end{tabular}

\vspace{2mm}
\parbox{0.95\textwidth}{\footnotesize Over ten seeds the classifier leads externally in 9 of 10 runs (Wilcoxon $p=0.014$) and internally in 7 of 10 ($p=0.041$). Relative to three seeds, the internal interval widens by a factor of about three and the external interval narrows by a factor of about four. Seeds 0--6 used batch 16 at learning rate $10^{-4}$ and seeds 7--9 batch 32 at $2\times10^{-4}$, a change made for compute reasons. The paired differences from the two groups are similar (internal $+2.44$ against $+3.09$; external $+2.87$ against $+2.61$; Welch $p=0.84$ and $p=0.92$), and we pool them on that basis. We state this as an assumption rather than as something a test established: with seven runs against three the Welch test has almost no power, and by the standard this paper applies elsewhere a non-significant result is not evidence of equivalence. Table~\ref{tab:perseed} prints every run so that a reader can inspect the two groups directly. Absolute accuracies differ slightly from Table~\ref{tab:h2h} because this replication uses one recipe for both arms rather than the per-arm settings of the original run; the paired differences, not the levels, are the quantity of interest.}
\label{tab:paired}
\end{table*}

\begin{table}[t]
\centering
\caption{Every run behind Table~\ref{tab:paired}. Accuracies are percentages; the difference columns are classifier minus pipeline, paired within a seed. Seeds 0--6 and 7--9 used different batch sizes and learning rates, so the two groups can be inspected separately.}
\footnotesize
\begin{tabular}{|c|c c c|c c c|}
\hline
 & \multicolumn{3}{c|}{Internal} & \multicolumn{3}{c|}{External} \\
Seed & Clf & Pipe & Diff & Clf & Pipe & Diff \\
\hline
0 & 84.9 & 82.9 & $+1.95$ & 79.6 & 74.9 & $+4.70$ \\
1 & 82.9 & 80.5 & $+2.44$ & 77.1 & 75.2 & $+1.88$ \\
2 & 84.4 & 79.0 & $+5.37$ & 77.1 & 75.9 & $+1.25$ \\
3 & 82.9 & 82.0 & $+0.98$ & 78.4 & 76.5 & $+1.88$ \\
4 & 84.9 & 77.6 & $+7.32$ & 79.3 & 73.4 & $+5.96$ \\
5 & 81.0 & 81.5 & $-0.49$ & 77.1 & 75.9 & $+1.25$ \\
6 & 83.9 & 84.4 & $-0.49$ & 80.3 & 77.1 & $+3.13$ \\
\hline
7 & 82.9 & 79.0 & $+3.90$ & 72.1 & 74.0 & $-1.88$ \\
8 & 81.5 & 83.4 & $-1.95$ & 78.1 & 73.4 & $+4.70$ \\
9 & 86.8 & 79.5 & $+7.32$ & 79.0 & 74.0 & $+5.02$ \\
\hline
Mean & 83.6 & 81.0 & $+2.63$ & 77.8 & 75.0 & $+2.79$ \\
\hline
\end{tabular}
\label{tab:perseed}
\end{table}

\subsection{Oracle versus real masks, and strict versus robust pipeline}
\label{Res:oracle}
Table~\ref{tab:h2h} also separates two effects that are easy to conflate. Swin-Tiny on ground-truth (oracle) masks reached 83.9 percent internally, an upper bound; the robust pipeline with real predicted masks reached 81.5 percent at the same seed and 78.9 $\pm$ 1.0 percent across seeds, so at matched seeds imperfect segmentation costs about two points. The choice of confidence threshold matters at the tail: the strict pipeline (confidence 0.25, no-detection counted as an error) reached 79.0 percent internally but only 54.2 percent externally, whereas the robust pipeline reached 69.9 percent externally at seed 42. We report both, and we never present the 83.9 percent oracle number as the accuracy of the deployed pipeline.

\subsection{External generalisation fails}
\label{Res:ext}
On genuinely clinical images the model does not transfer, and it fails in the dangerous direction. On the 94 BIP\_US images (Table~\ref{tab:ext}), the deployed pipeline graded about 55 of 94 images below their true severity and classified about half of the full-thickness burns as first degree; balanced accuracy was 30.5 percent under the deep-dermal-to-second-degree mapping and 35.0 percent under the deep-dermal-to-third-degree mapping, against a chance level of 50 percent for this two-class probe. On the clean external set of 319 images, the classifier was the most robust approach, with no undetected images, while the segmentation-dependent approaches degraded more (Table~\ref{tab:h2h}; Figure~\ref{fig:cm} shows the pipeline confusion matrices for both sets). Under-grading must be quantified conditionally, because a first-degree burn cannot be under-graded and a third-degree burn cannot be over-graded, so raw error counts are confounded by class balance. Among images that \emph{could} be under-graded, the pipeline under-graded 19.3 percent internally (23 of 119) and \textbf{28.3 percent externally (51 of 180)}; among images that could be over-graded it over-graded 10.1 percent internally (15 of 149) and 15.2 percent externally (45 of 296). The bias toward lower severity is therefore present on both sets and is not an artifact of the external class distribution. We had described the external rate as markedly higher; testing it rather than asserting it, the two-proportion comparison gives $z=1.77$, $p=0.08$, and that figure ignores the source clustering of Section~\ref{Mat:eval} and is therefore optimistic. The correct statement is that under-grading is substantial on both sets, numerically higher externally, and not established as different between them at this sample size. The perceptual-hash check that produced this set had already removed 118 training-identical images and 28 percent of sources before evaluation, so these numbers are measured on genuinely unseen data. Figure~\ref{fig:qual} shows what these errors look like in practice. The message is that the model does not generalise to clinical data, and that under-grading, directing a severe burn toward conservative care, is the more dangerous error.

\begin{table*}[t]
\centering
\caption{External validation of the deployed two-stage pipeline on the BIP\_US clinical database (94 images). Values are the number of images the pipeline predicted in each class, grouped by ground-truth burn depth. About 55 of 94 images were under-graded and about half of the full-thickness burns were called first degree.}
\small
\begin{tabular}{|l|c|c|c|c|}
\hline
Ground truth (depth) & Predicted first & Predicted second & Predicted third & $n$ \\
\hline
Superficial dermal & 13 & 27 & 2 & 42 \\
Deep dermal & 25 & 7 & 0 & 32 \\
Full thickness & 10 & 7 & 3 & 20 \\
\hline
\end{tabular}

\vspace{2mm}
\parbox{0.95\textwidth}{\footnotesize Balanced accuracy 30.5 percent under the deep-dermal-to-second-degree mapping and 35.0 percent under the deep-dermal-to-third-degree mapping, against a chance level of 50 percent for this two-class probe.}
\label{tab:ext}
\end{table*}

\begin{figure*}[t]
\centering
\includegraphics[width=\textwidth]{fig_qualitative.pdf}
\caption{Qualitative behaviour of the deployed two-stage pipeline on held-out internal test images, one per column: the input photograph, the region the localiser predicts, and the masked image that is passed to the classifier. The first three columns are correctly graded examples spanning the three severity classes. The fourth is a representative failure of the kind that matters clinically: a second-degree burn whose sparse, disconnected lesions are only partly recovered by the localiser, and which is graded first degree. Errors of this direction, toward lower severity, would route a patient toward conservative care and are the failure mode we quantify in Section~\ref{Res:ext}.}
\label{fig:qual}
\end{figure*}

\begin{figure*}[t]
\centering
\begin{minipage}{0.44\textwidth}
\centering
\includegraphics[width=\linewidth]{cm_true_pipeline.pdf}\\
(a)
\end{minipage}\hfill
\begin{minipage}{0.44\textwidth}
\centering
\includegraphics[width=\linewidth]{ext_cm_pipe.pdf}\\
(b)
\end{minipage}
\caption{Confusion matrices for the robust two-stage pipeline on a representative seed. Internally (a) errors fall mainly between neighbouring degrees, with second degree the hardest class, consistent with the segmentation result. On the clean external set (b) accuracy is lower and the more severe burns are frequently under-graded, for example second degree predicted as first, the dangerous direction discussed in Section~\ref{Res:ext}.}
\label{fig:cm}
\end{figure*}

\subsection{Any three of these ten seeds: the range of answers a three-seed protocol gives}
\label{Res:seeds}
Every multi-seed figure in the original protocol rests on three training runs, which is common practice but a thin basis for a variance estimate. The ten-seed replication of Section~\ref{Res:stat} lets us measure directly how thin (Table~\ref{tab:seeds}).

Over ten seeds the standalone classifier's internal accuracy has a sample standard deviation of \textbf{1.74} percentage points and a range of 5.85 points; the pipeline's is 2.22 points with a range of 6.83. Both are far larger than the three-seed protocol reported (0.75 and 1.02). The reason is visible directly: across all 120 possible three-seed subsets of these ten runs, the estimated standard deviation for the classifier ranges from \textbf{0.00 to 3.25} percentage points, and for the pipeline from 0.28 to 3.69. A three-seed standard deviation is therefore not a noisy estimate of the right quantity; it is very nearly uninformative, and it can be arbitrarily small purely by chance. Any interval built on one inherits that defect, which is exactly what Section~\ref{Res:stat} observed.

Section~\ref{Res:stat} compared three seeds against ten, but those two protocols also differ in training recipe, so that comparison confounds seed count with the recipe change. The confound is avoidable, because the ten runs already contain $\binom{10}{3}=120$ distinct three-seed subsets, and we can ask what each of them would have reported about exactly the same data. On the external comparison, where ten seeds give $+2.79$ percentage points with a 95 percent interval of $+1.09$ to $+4.49$ and $p=0.005$, the 120 subsets give point estimates from $+0.21$ to $+5.22$, $p$-values from $0.0005$ to $0.86$, and interval half-widths from $0.45$ to $10.63$ percentage points, a more than twenty-fold spread in apparent precision drawn from one set of runs. Only \textbf{15 of the 120} reach $p<0.05$; the remaining 105 would have reported no external difference. Internally the spread is wider still: point estimates run from $-0.98$ to $+6.67$ percentage points, so some three-seed subsets place the \emph{pipeline} ahead of the classifier, and only 6 of 120 reach $p<0.05$. Because every one of these subsets is drawn from the same ten runs under one recipe, this analysis carries no confound at all, and it isolates the seed count as the cause.

One observation cuts against us and we report it. Our original three-seed protocol reported an internal $p$ of 0.002, smaller than any of the 120 subsets produces (the minimum is $0.0094$). The original run and the replication differ in training recipe as well as in the runs themselves, so this is not a like-for-like comparison and we draw no inference from the gap. What it does establish is that an internal $p$ of 0.002 is not a stable property of this comparison, and we do not treat it as one.

We report this against our own results deliberately. It is the third of the six instances of the pattern this paper is about: a routine methodological shortcut, here too few seeds, manufactures apparent precision in the same way that splitting after augmentation manufactured an apparent effect. Section~\ref{Res:stat} shows what it costs in practice rather than in principle, and the practical recommendation is concrete: report the number of runs, prefer intervals on paired differences to error bars on individual systems, and budget for substantially more seeds than three.

\begin{table*}[t]
\centering
\caption{Seed variance measured over ten training runs on the internal test set, and how well three runs estimate it. The final column reports, across all 120 three-seed subsets of the ten runs, the smallest and largest standard deviation such a subset would have produced.}
\small
\begin{tabular}{|l|c|c|c|c|}
\hline
System & Seeds & Mean & Sample SD & SD from any 3 of these seeds \\
\hline
Standalone classifier (ConvNeXt-Large) & 10 & 83.61 & \textbf{1.74} & 0.00 to 3.25 \\
Two-stage pipeline (localiser to Swin-Tiny) & 10 & 80.98 & \textbf{2.22} & 0.28 to 3.69 \\
\hline
\end{tabular}

\vspace{2mm}
\parbox{0.95\textwidth}{\footnotesize For comparison, the three-seed standard deviations reported in Table~\ref{tab:h2h} are 0.75 and 1.02 percentage points, roughly \emph{half} the ten-seed values of 1.74 and 2.22. A three-seed estimate can be arbitrarily small by chance: for the classifier, some three-seed subsets of these ten runs give a standard deviation of zero.}
\label{tab:seeds}
\end{table*}

\begin{table}[t]
\centering
\caption{Inference efficiency on an NVIDIA A100 ($640\times640$ input, batch size 1). The full pipeline is about 1.6 times slower than the standalone classifier because it runs two models.}
\footnotesize
\begin{tabular}{|l|c|c|}
\hline
Stage or model & ms / image & Images/s \\
\hline
YOLO localiser (single class) & 15.4 & 65 \\
Standalone YOLOv8-seg (three class) & 15.1 & 66 \\
Swin-Tiny classifier & 23.7 & 42 \\
ConvNeXt-Large classifier & 25.8 & 39 \\
Full pipeline (localiser, mask, Swin) & 42.3 & 24 \\
\hline
\end{tabular}
\label{tab:eff}
\end{table}

\begin{table*}[t]
\centering
\caption{Skin-tone probe. Standalone classifier accuracy by Individual Typology Angle band, with Wilson 95 percent intervals, on both test sets. Bands are ordered lightest to darkest. Accuracy does not order monotonically with apparent skin tone on either set, and the band that is strongest internally is the weakest externally.}
\small
\begin{tabular}{|l|c|c|c|c|}
\hline
ITA band & \multicolumn{2}{c|}{Internal ($n=194$)} & \multicolumn{2}{c|}{External ($n=319$)} \\
 & $n$ & Accuracy (95\% CI) & $n$ & Accuracy (95\% CI) \\
\hline
Very light ($>55$) & 41 & 97.6 \ [87.4, 99.6] & 63 & 65.1 \ [52.8, 75.7] \\
Light (41 to 55) & 24 & 87.5 \ [69.0, 95.7] & 69 & 84.1 \ [73.7, 90.9] \\
Intermediate (28 to 41) & 40 & 72.5 \ [57.2, 83.9] & 58 & 63.8 \ [50.9, 74.9] \\
Tan (10 to 28) & 34 & 79.4 \ [63.2, 89.7] & 71 & 77.5 \ [66.5, 85.6] \\
Brown ($-30$ to 10) & 46 & 78.3 \ [64.4, 87.7] & 53 & 79.2 \ [66.5, 88.0] \\
Dark ($<-30$) & 9 & 77.8 \ [45.3, 93.7] & 5 & 80.0 \ [37.6, 96.4] \\
\hline
\end{tabular}

\vspace{2mm}
\parbox{0.95\textwidth}{\footnotesize Tested on the continuous ITA rather than on these bands, the classifier's errors fell on darker-appearing skin internally (median ITA 22.5 for errors against 34.2 for correct predictions, Mann-Whitney $p=0.015$) but not externally ($p=0.35$). The two-stage pipeline showed no association on either set ($p=0.42$ and $p=0.76$). Section~\ref{Res:ita} reports the subgroup analysis and its failure to replicate. The darkest band contains 9 and 5 images respectively; no claim about it is supportable, and we report it only so that the sparsity is visible.}
\label{tab:ita}
\end{table*}

\subsection{A skin-tone signal that did not replicate}
\label{Res:ita}
The probe of Section~\ref{Mat:ita} returned no evidence that either system's accuracy depends on apparent skin tone in a way that holds across datasets (Table~\ref{tab:ita}). Accuracy does not order monotonically with ITA on either test set, and the band that performs best internally, very light at 97.6 percent, is among the weakest externally at 65.1 percent, close to the external minimum of 63.8 percent in the intermediate band. That is the behaviour of noise rather than of a skin-tone effect.

We nonetheless report in full a subgroup result that we could easily have published as a positive finding, because how it failed is the point. Restricting the internal set to first-degree burns, the standalone classifier's errors fell on markedly darker-appearing skin than its correct predictions: median ITA \textbf{19.2} for the 21 errors against \textbf{38.4} for the 60 correct predictions, a 19-point gap that crosses a band boundary, from tan into intermediate, with Mann-Whitney $p=\mathbf{0.0025}$ and a rank-biserial correlation of $-0.45$. That is a large effect in the direction the fairness literature would predict and worry about, and it survives Bonferroni correction across all 15 tests in this probe (threshold $0.05/15=0.0033$). Reported alone, it would read as clear evidence that the classifier is biased against darker skin.

It does not replicate. The identical test on the external set, on the same arm and the same class, gives $p=\mathbf{0.35}$ with the effect in the \emph{opposite} direction (rank-biserial $+0.10$; median ITA 39.2 for errors against 37.6 for correct predictions). This is not a power failure: the external cell contains \textbf{139} images against the internal cell's 81, so the replication attempt was the better-powered of the two. No other cell in the probe approaches significance.

We draw the conservative conclusion, which is that this study provides \emph{no} usable evidence about skin-tone fairness in either direction, and that the internal subgroup result should be read as a false positive. We report it rather than delete it for the same reason we report the leakage artifact: this is the fourth time in this paper that a finding significant under one protocol has evaporated under a second look, and it is the most seductive of the six, because a bias result is one that a reader wants to believe and that a reviewer is unlikely to challenge. A single-dataset fairness analysis is subject to exactly the same replication failure as a single-dataset accuracy claim, and the fairness setting supplies extra incentive not to look twice. Establishing whether these models are equitable requires data with recorded phototype and a pre-registered analysis, neither of which we have.

\subsection{We fixed the leak in one stage and assumed the pipeline was clean}
\label{Res:segleak}
The four audits above were deliberate. This one was not: we found it while preparing the released code, after every other number in this paper was final, and it is the most uncomfortable result here because it is the paper's own subject matter applied to the paper's own pipeline.

Correcting the split (Section~\ref{Mat:leak}) fixed the classification datasets. It did not touch the segmentation models, which had been trained earlier on the collection's original split, and we did not think to check whether the two splits agreed. They do not: \textbf{179 of the 205 internal classification test images, 87.3 percent, lie in the segmentation models' training or validation folds}. Every internal number in this paper that passes through a YOLO model is therefore measured partly on images that model was fitted to.

We quantified the consequence rather than estimating it. Re-evaluating the contaminated three-class standalone model at the fair threshold of 0.05 gives \textbf{90.7 percent} on the internal test set (balanced 90.6) against \textbf{48.9 percent} externally (balanced 50.3). A 41.8-point internal-to-external collapse is not a distribution shift of that magnitude; it is the signature of a model being scored on its own training data.

We then measured the leak directly by removing it. We rebuilt both segmentation datasets on the classifiers' source-grouped split, so that the segmentation test fold \emph{is} the 205-image classification test fold and no source photograph appears in two folds, and retrained both YOLOv8x-seg models from the public initialisation under the original recipe: 100 epochs at $640\times640$, batch 8, AdamW at $10^{-3}$, seed 42, patience 20. Nothing changed except the partition. On the same 205 images, the three-class standalone model falls from 90.7 percent to \textbf{78.0 percent} (balanced 77.7, one image undetected), a paired difference that is highly significant on those images (exact McNemar, 33 against 7 discordant, $p=4\times10^{-5}$). The retrained single-class localiser reaches a mask mAP50 of 0.695 on that leak-free test fold and returns a region for all 205 images, so the robust pipeline's full-frame fallback never fires internally; the retrained three-class model reaches 0.566.

That 12.7-point gap is not purely the leak, and we would be repeating this paper's own error if we reported it as though it were. Correcting the split also shrank the training set, because the source-grouped partition assigns fewer images to training than the collection's original 90/6/4 split did: 2{,}425 against 3{,}081, a reduction of 21 percent. Some of the drop is therefore the loss of training data rather than the loss of an unfair advantage. The 26 test images that \emph{neither} model was trained on give a control for exactly this: on those, the contaminated model still scores 88.5 percent against the retrained model's 80.8, a gap of 7.7 points that cannot be leakage because neither model had seen them. On the 179 images the contaminated model had seen, the gap is 13.4 points. The difference between those two gaps, about 5.7 points, is the part attributable to leakage; but the control rests on 26 images and its bootstrap interval runs from 0.0 to 19.2 points, which is too wide to separate the two mechanisms cleanly. \textbf{What we can say is that removing the contamination cost 12.7 points overall, that a training-set reduction confounds it, and that a properly powered separation would need the leak-free model retrained at the original training-set size} --- which the source-grouped split makes impossible without collecting more photographs.

Two things follow. The standalone-YOLO row in Table~\ref{tab:h2h} can now be stated honestly at 78.0 percent, which places it below both the plain classifier and the pipeline and leaves the paper's ordering intact. And the internal pipeline arm, whose localiser carried the same contamination, is inflated by an amount we have now bounded empirically rather than guessed at, though not by the full 12.7 points, since only the pipeline's mask quality is affected and its classifier stage was always trained on the leak-free split.

Three things are \emph{not} affected, and we checked each rather than assuming it. The segmentation quality figures of Section~\ref{Res:seg} are computed on the segmentation models' own held-out split, which has zero source overlap with their training fold, so the mask mAP50 values of 0.726 and 0.603 stand. The external set has zero image-level and zero source-level overlap with the segmentation training data, so every external number in this paper is clean, including the entire ten-seed head-to-head on the external set. Benchmark 1 and the oracle-mask row use ground-truth polygon masks and never invoke a trained segmenter at all, so the paper's central masking null is untouched.

What it does affect is the \emph{internal} pipeline arm, whose localiser front-end produces better masks on these images than it would on unseen ones. That inflates the pipeline's internal accuracy. The direction matters: the pipeline is the arm that \emph{loses}, and it loses while enjoying this advantage. The internal margin of $+2.63$ percentage points in favour of the plain classifier is therefore a \emph{lower bound} on the true margin, and correcting the leak would widen it. The paper's conclusion survives, but it survives as a conservative statement rather than a clean one, and a reader is entitled to know the difference.

The transferable lesson is the one we most wish we had not had to learn from ourselves. A pipeline is only as leak-free as its leakiest stage, and correcting a split for one component does not correct it for a component trained earlier under a different regime. We had written a paper arguing that leakage is easy to introduce and hard to see, had built a source-grouped split specifically to eliminate it, had audited that split with perceptual hashing, and still shipped a two-stage system whose first stage had seen 87 percent of the test set. Nothing in our own procedure would have caught it; the check that did was comparing two directory listings.

\subsection{Efficiency}
\label{Res:eff}
The pipeline is slower and no more accurate than the standalone classifier (Table~\ref{tab:eff}). On the A100, the localiser took 15.4 ms per image and the standalone classifiers 25.8 ms (ConvNeXt-Large) and 23.7 ms (Swin-Tiny), while the full pipeline took 42.3 ms, about 1.6 times slower than the standalone classifier because it runs two models. The pipeline's return for this cost is localisation, not accuracy. These figures are measured on a datacentre A100 and are therefore not a mobile latency budget. We did not measure on-device inference time, and the shipped demonstration application runs a Swin-Small classifier rather than the benchmarked Swin-Tiny, so no number in this paper describes what a user of the application experiences.

\section{Discussion}
\label{Dis}
Six cautionary findings define this study, and each is an instance of the central theme: in every case a conventional protocol and a second, independent look gave different answers (Table~\ref{tab:audit}). First, leakage manufactured the masking benefit: augmenting a dataset before splitting it, then partitioning at the file level, produced a several-point improvement that does not exist. The failure is easy to make and hard to see, because the leaked copies are augmentations rather than exact duplicates. Second, identifier-based grouping, the standard correction for exactly that failure, still left 2 of 205 internal test images perceptually identical to a training image, one pair with conflicting labels, so grouping by identifier is a heuristic and not a proof of independence. Third, three training runs were too few to estimate either the head-to-head effect or its precision: the same ten runs that resolve the external margin at $p=0.005$ contain 105 of 120 three-seed subsets that would not have detected it at all, and 5 of 120 that place the pipeline ahead internally. Fourth, a single-dataset subgroup test handed us a large, Bonferroni-surviving skin-tone association that reversed direction and vanished on a better-powered external cell. Fifth, and least comfortably, our own corrected split covered only the classifiers: the segmentation models kept the collection's original split, and 87.3 percent of the internal test set sat inside their training or validation folds. A pipeline is only as leak-free as its leakiest stage. Sixth, the internal head-to-head margin that survives on plain accuracy does not survive on balanced accuracy, so which of two standard summaries we report decides whether that comparison is established at all; only the external margin holds under both.

Against that backdrop the burn-severity results are quickly stated: on a leak-free split masking gives no benefit resolvable at this sample size, the two-stage pipeline does not beat a plain classifier, and neither generalises to independent clinical images.

What makes the six instances worth reporting together is that they are not variations on a single mistake. They come from six different families of shortcut: a split procedure, a grouping heuristic, a seed count, a single-dataset subgroup test, an incompletely propagated correction, and the choice of summary metric, and in every case the first protocol was the ordinary one. The common structure is that a single evaluation cannot distinguish a finding from an artifact of the procedure that produced it, and that the researcher's own expectations determine which artifacts get a second look. We had no prior belief about masking or seeds, and we checked them. We did have a prior belief about skin tone, and we very nearly did not. And we believed our split was leak-free, so we did not re-audit it after correcting it, which is how the fifth one survived until the code was being packaged for release.

We nonetheless keep the segmentation-first design, and it is important to say why. Its justification is not accuracy, which we are careful not to claim, but localisation and interpretability: the system returns where the burn is and a clean masked region a clinician can inspect, which suits the mobile use-case and follows an established two-stage template \citep{bruno,almasni2020}. The honest reading is that any grading improvement segmentation could offer here is bounded well below the run-to-run variation of the training procedure, so what it reliably delivers is localisation, not accuracy.

The oracle-mask design strengthens rather than weakens this reading. Because the masking ablation used ground-truth masks, an upper bound, and still showed no benefit, poor segmentation cannot be blamed for the null: a better segmenter could not have rescued a masking gain that is absent even with perfect masks. The like-for-like gap at the same seed, 83.9 percent with oracle masks against 81.5 percent with real predicted masks, about two points, is the cost of imperfect masks, not a hidden pipeline accuracy.

The high single-dataset accuracies reported in the burn literature, 95.6 percent for a transfer-learning classifier \citep{suha} and up to 98.1 percent among the studies a recent review collates \citep{fu2026review}, are of the kind our own leaky configurations produced before correction, which is exactly why such numbers deserve caution when the dataset is augmented and the split procedure is not described. The mirror image is instructive: multi-source training improves robustness \citep{elsarta2025}, whereas our single-source model fails to transfer. Our external failure is consistent with a wide medical-imaging pattern in which strong internal scores collapse on external data \citep{zech2018,yu2022,martensson2020,roberts}.

External behaviour is at once the strongest limitation and the strongest message, and it differs by set. On the genuinely clinical Seville images the model collapses toward chance and under-grades severely. On the second, web-sourced external set the decline is more modest and the plain classifier remains the most robust approach, consistent with a smaller distribution shift for web images than for clinical photographs. In neither case is the system ready for a clinical role. The clinical collapse most likely reflects a narrow, web-sourced training distribution and a labelling protocol that differs from clinical depth staging. Multi-source, clinically labelled data and domain-adaptation methods are needed before generalisation can be claimed.

\subsection{Limitations}
\label{Dis:lim}
Several limitations qualify these results. The most serious is the one in Section~\ref{Res:segleak}: our source-grouped split was built for the classifiers and never propagated to the segmentation models, so 87.3 percent of the internal test set lies in their training or validation folds. We retrained both segmentation models on the source-grouped split, which fixes the standalone-YOLO row, but we did not re-run the two-stage pipeline arm, and we state the reason accurately rather than flatteringly. The pipeline's Swin-Tiny classifier is trained on ground-truth masks and meets the localiser only at inference, so substituting the retrained localiser is an inference-only operation; the masks for all 205 test images already exist; and the ten-seed recipe is documented in this paper (Table~\ref{tab:paired}). The obstacle is therefore neither conceptual nor a missing recipe, as an earlier draft of this paragraph asserted: it is simply that the ten Swin-Tiny checkpoints were not archived, so closing this would mean retraining that arm, which we have not done. The internal pipeline accuracies in Table~\ref{tab:h2h} consequently still carry the original localiser and remain optimistic, the classifier's internal margin over the pipeline is a lower bound rather than a point estimate, and that comparison should be read as directional. We flag this as the most valuable single experiment a reader could run against our released artifacts. The external comparison, the segmentation quality figures and the whole of Benchmark 1 are unaffected, and we verified each rather than assuming it. The training data come from a single web source with community labels that were not clinically verified. The internal test set is small ($N=205$, one image per source), so confidence intervals are wide. The head-to-head standalone YOLOv8-seg model is single-seed, so its point estimate should not be compared to the multi-seed means without that caveat. The masking delta is not stable to the architecture pool or to run-to-run noise, so we make no claim that masking is exactly zero. The headline multi-seed figures rest on three training runs, and Section~\ref{Res:seeds} shows directly that three runs estimate the seed variance poorly: a ten-seed replication of the classifier gave a standard deviation of 1.74 percentage points against 0.75 from three, and any three of those ten would have yielded anywhere from 0.00 to 3.25. Repeating the full head-to-head over ten seeds (Table~\ref{tab:paired}) changed which comparison the evidence most strongly supports, so the three-seed intervals in Table~\ref{tab:h2h} should not be read as reliable; we retain them only because they are the original protocol. Future work of this kind should budget for substantially more seeds than is conventional. The BIP\_US external set is small and lacks first-degree cases, so it probes second versus third degree only. We also conducted no user study or prospective clinical evaluation of the mobile application, so its usability and clinical value remain untested. Finally, our skin-tone analysis is a proxy measurement and must not be read as a fairness audit: ITA estimated from uncontrolled photographs tracks apparent skin tone in the image rather than constitutive phototype, the darkest band holds only 9 and 5 images, and the one subgroup signal we found did not replicate (Section~\ref{Res:ita}). Whether these models are equitable across skin tones remains genuinely unknown, and answering it needs data with recorded phototype and a pre-registered analysis.

\section{Conclusion}
\label{Con}
We studied a deployed two-stage burn-severity system through a source-grouped, externally validated benchmark, and six routine evaluation choices changed the answer. Splitting an augmented dataset at the file level manufactured a masking benefit that vanished under a source-grouped split. Grouping by source identifier still left 1.0 percent of the internal test set perceptually identical to training, one pair with conflicting labels. A three-seed protocol overstated the internal head-to-head margin, narrowed its interval by a factor of about three, and concealed the external margin entirely; any three of our ten replication runs would have reported a classifier standard deviation somewhere between 0.00 and 3.25 percentage points, and a pipeline one between 0.28 and 3.69. A single-dataset subgroup test produced a large, Bonferroni-surviving skin-tone association that reversed and vanished externally. And our corrected split reached only the classifiers: the segmentation stage kept the original split, putting 87.3 percent of the internal test set inside its training or validation folds. Retraining it on the source-grouped split cost the standalone segmentation model 12.7 percentage points on those same images, though that figure is confounded with the smaller training set the correction implies; externally, where there was never any contamination to remove, the retrained model is 17.9 points better. Under the corrected protocol a plain classifier was the most accurate single system and beat the two-stage pipeline on external data, which is uncontaminated, in nine of ten seeds, so segmentation-first is justified here by localisation and interpretability rather than by accuracy; and the model did not transfer to clinical images and under-graded severity. The procedural lessons transfer beyond burns: datasets augmented before splitting must be partitioned by source image, or reported gains may be illusory; identifier grouping should be verified with a content-based check; three seeds neither estimate a variance nor support an interval; a single-dataset subgroup finding, a fairness finding included, is subject to the same replication failure as a single-dataset accuracy claim; and a correction applied to one stage of a pipeline must be propagated to every stage and then re-audited, because a pipeline is only as leak-free as its leakiest component. In-distribution accuracy must not be read as clinical usefulness. Future work needs multi-source, clinically labelled data with patient-level splits and domain adaptation, and should keep segmentation for localisation rather than as an accuracy lever.

% =====================================================================
%  MELBA mandatory closing sections
% =====================================================================

\acks{%
This study began as the authors' undergraduate graduation project at the Faculty of Computer and Information Systems, Islamic University of Madinah. All three authors have since graduated; the additional work reported here, comprising the leakage audit, the leak-free re-analysis, the ten-seed replication and the external validation, was carried out independently, without institutional supervision, funding or compute. This research received no external funding.

The authors thank their supervisor Dr.\ Hani Sayaf Almoamari for guidance throughout the original graduation project, the Biomedical Image Processing Group of the University of Seville for making the BIP\_US database available for research, and the Google Colab and Kaggle platforms for computational resources.

\textbf{Declaration of generative AI and AI-assisted technologies.} During the preparation of this work the authors used Anthropic Claude (Claude Code; Opus 4.8 and Opus 5 models) to organise the project materials, reproduce and check the reported results, run the statistical, benchmarking and external-validation analyses (including the leak-free re-runs, the ten-seed replication, the perceptual-hash leakage check and the skin-tone probe), and to draft and edit the manuscript text. All drafting was performed from author-supplied results and under author direction; the authors verified every quantitative claim against the underlying result files, and the repository cited in the Data Availability Statement contains a script that recomputes each published number from the raw data. No text was accepted without author review, and the tool was not used for unsupervised, de novo content creation. The research itself, comprising the system design, model training, dataset preparation and the mobile application, was conceived and carried out by the authors, who take full responsibility for the content of this publication.}

\ethics{%
This study used only publicly available, de-identified burn images together with the BIP\_US research database, which was obtained from the Biomedical Image Processing Group of the University of Seville for research use. No new data were collected from human participants and no identifiable individuals appear in the reported figures. Institutional review board approval was therefore not required. The work follows appropriate ethical standards in conducting the research and in writing the manuscript. We note explicitly that the system described here is a research and demonstration artifact; the paper reports that it does not generalise to independent clinical images, and it must not be used for clinical decision-making.}

\coi{%
We declare we don't have conflicts of interest. The authors have no financial, organizational, commercial or personal conflicts of interest to disclose, and have not recently collaborated with any member of the MELBA editorial board.}

\data{%
All data supporting the findings of this study are available. The primary burn dataset is the publicly available ``skin burn wound classification'' collection on Roboflow Universe \citep{roboflow}, released under a Creative Commons Attribution 4.0 licence. The BIP\_US clinical database is available for research from the Biomedical Image Processing Group of the University of Seville, on request to its custodians; we redistribute no BIP\_US image. Code, the leak-free split definition and its builder, the raw per-image predictions behind every table and figure, the perceptual-hash contamination lists, and the Kaggle notebooks that produced the segmentation retraining of Section~\ref{Res:segleak} are available at \url{https://github.com/Just-Ryan/burn-severity-benchmark}. Trained weights exceed the repository host's per-file limit and are attached to the tagged release there: both segmentation models retrained on the source-grouped split, both originals trained on the collection's distribution-supplied split, and the classifier shipped in the mobile demonstration. We publish the contaminated models alongside the corrected ones deliberately, so that the leak of Section~\ref{Res:segleak} can be reproduced as readily as its correction. The benchmark's Swin-Tiny classifier weights were not archived at the time and are not available; this is why the two-stage pipeline arm could not be re-evaluated on the corrected localiser. That repository includes \texttt{verify\_paper\_numbers.py}, a self-contained script that recomputes \textbf{156} published quantities from the committed raw data without a GPU, a dataset download or model weights, and reports each against the value printed here; all 156 pass. Two of the checks exist specifically to catch us overstating our own results, one confirming that the matched-architecture leakage contrast is the smaller $2.07 \rightarrow 0.55$ rather than $3.06 \rightarrow 0.55$, and one confirming that the internal interval-narrowing factor is 3.35 rather than the ``four times'' an earlier draft of this manuscript claimed. The repository also includes \texttt{skin\_tone\_probe.py}, which reproduces the ITA analysis of Section~\ref{Res:ita} from the test images and the released per-image predictions.}

% ---------------------------------------------------------------------
% melba.cls defines \doi{} as a *setter* for this paper's own DOI header
% field: \def\doi#1{\gdef\@doi{#1}} — it stores the argument and prints
% nothing. plainnat.bst emits \doi{...} for every reference, so without
% this fix (a) every reference DOI is silently swallowed, leaving a bare
% full stop, and (b) the last reference's DOI overwrites \@doi, so an
% accepted paper would print someone else's DOI in its own header and
% Crossmark link. The header is typeset by \maketitle above, so \@doi has
% already been consumed and it is safe to repurpose \doi from here on.
% ---------------------------------------------------------------------
\makeatletter
\let\melba@headerdoi\doi
\renewcommand{\doi}[1]{\href{https://doi.org/#1}{doi:#1}}
\makeatother

\bibliography{melba}

\end{document}
